\chapter{The Central Limit Theorem, Sharpened}
\begin{center}\itshape
Volume~4 proved that sums of many small independent pieces look Gaussian. This chapter asks the questions that come next:\\
is the condition of the theorem the right one? How good is the Gaussian approximation for a given $n$?\\
Does it hold point by point, for the probability of each value? And what does it say about walks coming home?
\end{center}

In 1920 George P\'olya wrote a paper on sums of random variables and the problem of moments, and titled it ``On the central limit theorem of probability
theory''. The theorem had no name before; P\'olya called it \emph{central} because of the central place it holds in probability, and the name stuck,
adjective and all. At that time the theorem had been proved under several sets of conditions (by Lyapunov in 1900 and 1901, assuming moments of order a little above two), but nobody
knew where the truth lay.

Two years later the Finnish mathematician Jarl Waldemar Lindeberg found both a short new proof, the replacement method of Volume~4, and the condition
that now carries his name. Was it the right condition? The answer came in 1935, from two people at once. William Feller, a young refugee then in
Stockholm, and Paul L\'evy in Paris proved, independently, that Lindeberg's condition is also \emph{necessary}, as soon as no single summand
matters. Lucien Le Cam later wrote the history of that year, and of who knew what when, in a well-known essay.

Then came precision. In 1941 Andrew Berry in the United States and in 1942 Carl-Gustav Esseen in Sweden, separately, showed that the error of the
Gaussian approximation is at most a constant times $\rho/(\sigma^3\sqrt n)$. The constant has a story of its own: Esseen's first value was $7.59$; it has
been pushed down over seventy years to $0.4748$ (Shevtsova, 2011), and Esseen himself showed in 1956 that it cannot be below $0.4097$. The gap is
still open.

The chapter proves these four results in that order: Lindeberg--Feller, Berry--Esseen, the multivariate theorem with Pearson's $\chi^2$ test as
a test case, and the local limit theorem. The last section turns the Fourier method on an old friend: when does a random walk on $\Z^d$ come back home?

\begin{reminder}[from Volume 4 and Chapter 3]
The \emph{characteristic function} of a random vector $X$ in $\R^d$ is $\varphi_X(t)=\mathbb E e^{i\langle t,X\rangle}$; it determines the law, and if
$\int|\varphi_X|<\infty$ for real $X$ then $X$ has the density $\frac1{2\pi}\int e^{-itx}\varphi_X(t)\dd t$ (Volume~4, Theorem~6.8). The \emph{Lindeberg
theorem} (Volume~4, Theorem~6.10): if $X_{n,1},\dots,X_{n,k_n}$ are independent and centred, with variances $\sigma_{n,k}^2$ summing to $1$, and for
every $\varepsilon>0$
\[
L_n(\varepsilon)=\sum_k\mathbb E\bigl[X_{n,k}^2\mathbf 1_{|X_{n,k}|>\varepsilon}\bigr]\to0,
\]
then $S_n=\sum_kX_{n,k}\Rightarrow N(0,1)$. In particular (Volume~4, Corollary~6.11) $S_n/(\sigma\sqrt n)\Rightarrow N(0,1)$ for i.i.d.\ centred steps of
variance $\sigma^2$. \emph{L\'evy's continuity theorem} (Theorem~\ref{thm:levycont}): $X_n\Rightarrow X$ in $\R^d$ if and only if
$\varphi_{X_n}\to\varphi_X$ pointwise. The \emph{Cram\'er--Wold device} (Corollary~\ref{cor:cramerwold}): $X_n\Rightarrow X$ in $\R^d$ if and
only if $\langle t,X_n\rangle\Rightarrow\langle t,X\rangle$ for every $t$. For a walk on $\Z^d$ with i.i.d.\ steps, $u_n=P(S_n=0)$, and the walk is
\emph{recurrent} (returns to $0$ with probability $1$, then infinitely often) if and only if $\sum_nu_n=\infty$ (Volume~4, Theorem~7.12: the proof
there is written for the simple walk but uses only that the steps are i.i.d.).
\end{reminder}

We use two Taylor bounds for $u$ real, both from $e^{iu}-\sum_{j<m}\frac{(iu)^j}{j!}=\frac{i^m}{(m-1)!}\int_0^u(u-s)^{m-1}e^{is}\dd s$:
\begin{equation}\label{eq:taylor}
\Bigl|e^{iu}-1-iu\Bigr|\le\frac{u^2}2,\qquad\Bigl|e^{iu}-1-iu+\frac{u^2}2\Bigr|\le\frac{|u|^3}6.
\end{equation}

\section{Lindeberg's condition is the right one}

Lindeberg's condition says that the summands which are large, compared with the total standard deviation, carry a negligible part of the variance.
It implies that every single summand is negligible: $\sigma_{n,k}^2\le\varepsilon^2+L_n(\varepsilon)$, so
\begin{equation}\label{eq:uan}
\max_k\sigma_{n,k}^2\to0.
\end{equation}
Without~\eqref{eq:uan} there is no hope of a converse: if $X_{n,1}\sim N(0,1)$ and all the other summands are $0$, the sum is Gaussian and Lindeberg's
condition fails. Feller's theorem says that this is the only way it can fail.

\begin{theorem}[Lindeberg--Feller]\label{thm:feller}
Let $X_{n,1},\dots,X_{n,k_n}$ be independent and centred, with variances $\sigma_{n,k}^2$ such that $\sum_k\sigma_{n,k}^2=1$ and $\max_k\sigma_{n,k}^2\to0$.
Then $S_n\Rightarrow N(0,1)$ if and only if Lindeberg's condition holds.
\end{theorem}
\begin{proof}
``If'' is Volume~4's theorem. Conversely, suppose $S_n\Rightarrow N(0,1)$, and fix $t\in\R$. Let $\varphi_{n,k}$ be the characteristic function of
$X_{n,k}$. By~\eqref{eq:taylor} and $\mathbb EX_{n,k}=0$, $|\varphi_{n,k}(t)-1|\le\frac{t^2}2\sigma_{n,k}^2$, which is at most $\frac12$ for all $k$ once $n$ is
large. For complex $w$ with $|w|\le\frac12$, the series of the logarithm gives
\[
|\log(1+w)-w|\le\sum_{j\ge2}\frac{|w|^j}2=\frac{|w|^2}{2(1-|w|)}\le|w|^2.
\]
Taking real parts, with $w=\varphi_{n,k}(t)-1$ and $\operatorname{Re}\log z=\log|z|$:
\[
\Bigl|\sum_k\log|\varphi_{n,k}(t)|-\sum_k\operatorname{Re}(\varphi_{n,k}(t)-1)\Bigr|\le\sum_k|\varphi_{n,k}(t)-1|^2\le\frac{t^2}2\max_k\sigma_{n,k}^2\cdot\frac{t^2}2\sum_k\sigma_{n,k}^2\to0.
\]
The first sum is $\log|\varphi_{S_n}(t)|$, which tends to $\log e^{-t^2/2}=-\frac{t^2}2$ because the characteristic functions converge. Hence
\[
\sum_k\mathbb E\bigl[1-\cos(tX_{n,k})\bigr]=-\sum_k\operatorname{Re}(\varphi_{n,k}(t)-1)\to\frac{t^2}2.
\]
Now split each expectation at $|X_{n,k}|=\varepsilon$. On $\{|X_{n,k}|\le\varepsilon\}$ use $1-\cos u\le\frac{u^2}2$; on the complement use $1-\cos u\le2$ and
Chebyshev, $\sum_kP(|X_{n,k}|>\varepsilon)\le\sum_k\sigma_{n,k}^2/\varepsilon^2=\varepsilon^{-2}$:
\[
\sum_k\mathbb E\bigl[1-\cos(tX_{n,k})\bigr]\le\frac{t^2}2\sum_k\mathbb E\bigl[X_{n,k}^2\mathbf 1_{|X_{n,k}|\le\varepsilon}\bigr]+\frac2{\varepsilon^2}
=\frac{t^2}2\bigl(1-L_n(\varepsilon)\bigr)+\frac2{\varepsilon^2}.
\]
Letting $n\to\infty$, $\frac{t^2}2\limsup_nL_n(\varepsilon)\le\frac2{\varepsilon^2}$. This holds for every $t$, so $\limsup_nL_n(\varepsilon)=0$.
\end{proof}

The proof shows where the Gaussian comes from: $\sum_k\mathbb E[1-\cos tX_{n,k}]$ must behave like $\frac{t^2}2$, the contribution of small values,
and large values would add a part bounded in $t$, which a quadratic cannot absorb. When the large values are \emph{not} negligible the limit can
still exist, but it is not Gaussian: that is the subject of Chapter~5 (Poisson and stable limits).

\begin{corollary}[Lyapunov's condition]\label{cor:lyapunov}
Under the normalisation of Theorem~\ref{thm:feller}, if $\sum_k\mathbb E|X_{n,k}|^{2+\delta}\to0$ for some $\delta>0$, then $S_n\Rightarrow N(0,1)$.
\end{corollary}
\begin{proof}
On $\{|X|>\varepsilon\}$, $X^2\le|X|^{2+\delta}/\varepsilon^\delta$. So
\[
L_n(\varepsilon)\le\varepsilon^{-\delta}\sum_k\mathbb E|X_{n,k}|^{2+\delta}\to0.
\]
\end{proof}

\begin{example}[inversions, or the work of insertion sort]\label{ex:inversions}
An \emph{inversion} of a permutation $\pi$ of $\{1,\dots,n\}$ is a pair of positions $i<j$ with $\pi(i)>\pi(j)$. Insertion sort, which slides each new
element left past the larger ones, makes exactly one swap per inversion, so the number $I_n$ of inversions of a uniform random permutation is the
cost of insertion sort on random data. For $k\le n$ let $Y_k$ be the number of $i<k$ with $\pi(i)>\pi(k)$: the number of earlier elements that the
$k$-th must pass. Then $I_n=Y_1+\dots+Y_n$ and $0\le Y_k\le k-1$. The map $\pi\mapsto(Y_1,\dots,Y_n)$ is injective (read the vector backwards:
$Y_n$ says which value sits in the last position, namely the $(Y_n+1)$-th largest, and so on), and both sets have $n!$ elements, so it is a
bijection onto $\prod_k\{0,\dots,k-1\}$. Under the uniform law, therefore, the $Y_k$ are \emph{independent}, $Y_k$ uniform on $\{0,\dots,k-1\}$, with
mean $\frac{k-1}2$ and variance $\frac{k^2-1}{12}$. Summing (Exercise~\ref{ex:inversionsvar}),
\[
\mathbb EI_n=\frac{n(n-1)}4,\qquad s_n^2=\operatorname{Var}I_n=\frac{2n^3+3n^2-5n}{72}.
\]
Put $X_{n,k}=(Y_k-\frac{k-1}2)/s_n$. Then $|X_{n,k}|\le n/s_n\to0$, so for each $\varepsilon$ all the indicators in $L_n(\varepsilon)$ vanish for
large $n$, and Lindeberg gives
\[
\frac{I_n-n^2/4}{n^{3/2}/6}\Rightarrow N(0,1).
\]
Sorting a million random keys by insertion costs $2.5\times10^{11}$ swaps, give or take about $1.7\times10^8$: the fluctuation is less than one part in a
thousand.
\end{example}

\section{How fast? The Berry--Esseen theorem}

The central limit theorem says $F_n(x)\to\Phi(x)$ for every $x$ (Volume~4, Proposition~6.3). We now bound the speed, uniformly in $x$. The tool is a smoothing inequality: it converts a bound on characteristic functions into a bound on distribution functions,
at the price of a term that measures how much we smoothed.

\begin{lemma}[P\'olya's density]\label{lem:polya}
$h(x)=\dfrac{1-\cos x}{\pi x^2}$ is a probability density, with characteristic function $(1-|t|)^+$. For $T>0$, $h_T(x)=Th(Tx)$ is a density
with characteristic function $(1-|t|/T)^+$, and if $Y$ has density $h_T$ then $P(|Y|>\delta)\le\frac4{\pi T\delta}$.
\end{lemma}
\begin{proof}
The triangular density $(1-|y|)^+$ has characteristic function $2\int_0^1(1-y)\cos(ty)\dd y=\frac{2(1-\cos t)}{t^2}=2\pi h(t)$, which is integrable. By
inversion, $(1-|y|)^+=\frac1{2\pi}\int e^{-ity}\,2\pi h(t)\dd t=\int e^{-ity}h(t)\dd t$. At $y=0$ this says $\int h=1$, and $h\ge0$; replacing $y$ by $-t$,
the characteristic function of $h$ is $(1-|t|)^+$. Scaling gives $h_T$. Finally $h_T(x)\le\frac2{\pi Tx^2}$, so
$P(|Y|>\delta)\le2\int_\delta^\infty\frac2{\pi Tx^2}\dd x=\frac4{\pi T\delta}$.
\end{proof}

\begin{lemma}[Esseen's smoothing inequality]\label{lem:smoothing}
Let $F$ and $G$ be distribution functions with characteristic functions $\varphi$ and $\gamma$, where $G$ has a density bounded by $m$. Let $T>0$, and
suppose $\int_{-T}^T|\varphi(t)-\gamma(t)|/|t|\dd t<\infty$. Then
\[
\sup_x|F(x)-G(x)|\le\frac1\pi\int_{-T}^T\frac{|\varphi(t)-\gamma(t)|}{|t|}\dd t+\frac{24m}{\pi T}.
\]
\end{lemma}
\begin{proof}
Let $\Delta=F-G$, $\eta=\sup|\Delta|$, and $\Delta_T(x)=\int\Delta(x-y)h_T(y)\dd y$.

\emph{Step 1: $|\Delta_T(x)|\le\frac1{2\pi}\int_{-T}^T|\varphi-\gamma|/|t|$.} The function $\int F(x-y)h_T(y)\dd y$ is the distribution function of
$X+Y$, with $X\sim F$ and $Y\sim h_T$ independent. Its characteristic function $\varphi(t)(1-|t|/T)^+$ is integrable, so it has the density
$\frac1{2\pi}\int_{-T}^Te^{-itx}\varphi(t)(1-|t|/T)\dd t$, and the same holds for $G$. Hence
\[
\Delta_T'(x)=\frac1{2\pi}\int_{-T}^Te^{-itx}\bigl(\varphi(t)-\gamma(t)\bigr)\Bigl(1-\frac{|t|}T\Bigr)\dd t.
\]
Let $R(x)=\frac1{2\pi}\int_{-T}^Te^{-itx}g(t)\dd t$ with $g(t)=\frac{\varphi(t)-\gamma(t)}{-it}(1-\frac{|t|}T)$, integrable by hypothesis. Differentiating
under the integral (Volume~4, Lemma~6.6, with the bound $|\varphi-\gamma|$), $R'=\Delta_T'$. Both functions tend to $0$ at $-\infty$: $\Delta_T$ by
dominated convergence, and $R$ because $\int e^{-itx}g(t)\dd t\to0$ as $|x|\to\infty$ for integrable $g$ (true for indicators of intervals by direct
computation, hence for step functions, which are dense in $L^1$). So $\Delta_T=R$, and $|R|\le\frac1{2\pi}\int|g|$.

\emph{Step 2: $\eta\le2\sup|\Delta_T|+\frac{24m}{\pi T}$.} Assume $\eta>0$. Since $\Delta\to0$ at $\pm\infty$ and $\Delta$ has one-sided limits, there is $x_0$
such that $\Delta(x_0)$ or $\Delta(x_0-)$ equals $\eta$ or $-\eta$. Suppose it is $+\eta$. As $F$ is non-decreasing and $G$ is $m$-Lipschitz, for $s>0$,
$\Delta(x_0+s)\ge\eta-ms$. Let $\delta=\frac\eta{2m}$ and $x_1=x_0+\delta$. For $|y|<\delta$, $x_1-y=x_0+(\delta-y)$ with $\delta-y>0$, so
$\Delta(x_1-y)\ge\eta-m(\delta-y)=\frac\eta2+my$. Using $\Delta\ge-\eta$ elsewhere, the symmetry of $h_T$, and Lemma~\ref{lem:polya},
\[
\Delta_T(x_1)\ge\int_{|y|<\delta}\Bigl(\frac\eta2+my\Bigr)h_T(y)\dd y-\eta P(|Y|\ge\delta)\ge\frac\eta2-\frac{3\eta}2\cdot\frac4{\pi T\delta}=\frac\eta2-\frac{12m}{\pi T}.
\]
If it is $-\eta$, the mirror argument ($\Delta(x_0-s)\le-\eta+ms$ for $s>0$, and $x_1=x_0-\delta$) gives $-\Delta_T(x_1)\ge\frac\eta2-\frac{12m}{\pi T}$.
Steps 1 and 2 together give the lemma.
\end{proof}

\begin{lemma}\label{lem:cfestimate}
Let $X$ be centred with $\mathbb EX^2=1$ and $\rho=\mathbb E|X|^3<\infty$, with characteristic function $\varphi$. Then $\rho\ge1$, and for $n\ge10$ and
$|t|\le T=\frac{4\sqrt n}{3\rho}$,
\[
\bigl|\varphi(t/\sqrt n)^n-e^{-t^2/2}\bigr|\le\frac{\rho|t|^3}{3\sqrt n}e^{-t^2/4}.
\]
\end{lemma}
\begin{proof}
$\rho\ge(\mathbb EX^2)^{3/2}=1$ by Jensen (or H\"older). Let $s=t/\sqrt n$, so $|s|\le\frac4{3\rho}\le\frac43$. By~\eqref{eq:taylor}, $|\varphi(s)-1+\frac{s^2}2|\le\frac{\rho|s|^3}6$.
As $\frac{s^2}2\le\frac89<1$ and $\frac{\rho|s|}6\le\frac29$,
\[
|\varphi(s)|\le1-\frac{s^2}2+\frac{\rho|s|}6s^2\le1-\frac5{18}s^2\le e^{-5s^2/18}.
\]
So $|\varphi(s)|^{n-1}\le e^{-5(n-1)t^2/(18n)}\le e^{-t^2/4}$, since $20(n-1)\ge18n$ for $n\ge10$; also $(e^{-s^2/2})^{n-1}\le e^{-t^2/4}$. For complex $a,b$ of
modulus at most $1$, $|a^n-b^n|=|a-b|\,|\sum_ja^jb^{n-1-j}|\le n|a-b|\max(|a|,|b|)^{n-1}$. With $a=\varphi(s)$, $b=e^{-s^2/2}$, and $0\le e^{-v}-1+v\le\frac{v^2}2$
for $v\ge0$,
\[
n|a-b|\le n\Bigl(\frac{\rho|s|^3}6+\frac{s^4}8\Bigr)=\frac{\rho|t|^3}{6\sqrt n}+\frac{t^4}{8n}.
\]
Finally $\frac{t^4}{8n}\le\frac{|t|^3}{8n}\cdot\frac{4\sqrt n}{3\rho}=\frac{|t|^3}{6\rho\sqrt n}\le\frac{\rho|t|^3}{6\sqrt n}$, since $\rho\ge1$.
\end{proof}

\begin{theorem}[Berry--Esseen]\label{thm:berryesseen}
Let $X_1,X_2,\dots$ be i.i.d.\ and centred, with variance $\sigma^2>0$ and $\rho=\mathbb E|X_1|^3<\infty$, and let $F_n$ be the distribution function of
$S_n/(\sigma\sqrt n)$. Then for every $n\ge1$
\[
\sup_{x\in\R}|F_n(x)-\Phi(x)|\le\frac{3.1\,\rho}{\sigma^3\sqrt n}.
\]
\end{theorem}
\begin{proof}
Replacing $X_i$ by $X_i/\sigma$, we may take $\sigma=1$. If $n\le9$, the right side is at least $3.1/3>1$ (as $\rho\ge1$) and there is nothing to prove.
For $n\ge10$, apply Lemma~\ref{lem:smoothing} with $F=F_n$, $G=\Phi$, $m=\frac1{\sqrt{2\pi}}$ and $T=\frac{4\sqrt n}{3\rho}$. By Lemma~\ref{lem:cfestimate} the
integrand is at most $\frac{\rho t^2}{3\sqrt n}e^{-t^2/4}$, which is integrable, and $\int_\R t^2e^{-t^2/4}\dd t=4\sqrt\pi$. So
\[
\sup_x|F_n(x)-\Phi(x)|\le\frac1\pi\cdot\frac{4\sqrt\pi\,\rho}{3\sqrt n}+\frac{24}{\pi\sqrt{2\pi}}\cdot\frac{3\rho}{4\sqrt n}=\Bigl(\frac4{3\sqrt\pi}+\frac{18}{\pi\sqrt{2\pi}}\Bigr)\frac\rho{\sqrt n},
\]
and the bracket is $0.7523\ldots+2.2858\ldots<3.1$.
\end{proof}

The rate $n^{-1/2}$ cannot be improved in general, and the reason is arithmetic: a lattice variable puts atoms of size about $n^{-1/2}$, which no
continuous law can match.

\begin{proposition}\label{prop:latticelower}
For the coin-tossing walk ($X_i=\pm1$, so $\sigma=\rho=1$) and $n$ even, $\sup_x|F_n(x)-\Phi(x)|\ge\dfrac1{2\sqrt{2n}}$.
\end{proposition}
\begin{proof}
$F_n$ jumps by $u_n=P(S_n=0)$ at $0$, while $\Phi$ is continuous, so $|F_n(0)-\Phi(0)|+|\Phi(0)-F_n(0-)|\ge u_n$ and one of the two terms is at least
$\frac{u_n}2$. By Volume~4, Lemma~7.1, $u_n\ge\frac1{2\sqrt{n/2}}=\frac1{\sqrt{2n}}$.
\end{proof}

\begin{remark}[the constant]
Our $3.1$ is far from the truth. The best constant $C_0$ in $\sup|F_n-\Phi|\le C_0\rho/(\sigma^3\sqrt n)$ is known to satisfy
\[
0.4097\le\frac{\sqrt{10}+3}{6\sqrt{2\pi}}\le C_0\le0.4748;
\]
the lower bound is Esseen's (1956), attained in the limit by suitable two-point laws, and the upper bound is Shevtsova's (2011), the latest in a long series of
improvements. Proposition~\ref{prop:latticelower} gives $C_0\ge0.35$ for free. Figure~\ref{fig:berryesseen} shows the exact error for binomial sums.
\end{remark}

\begin{figure}[ht]
\centering\fig{ch04-berryesseen}
\caption{The exact Kolmogorov distance $D_n=\sup_x|F_n(x)-\Phi(x)|$ for sums of $n$ Bernoulli$(p)$ variables, multiplied by $\sqrt n$ and divided by
$\rho/\sigma^3$, for $p=\frac12$ (black) and $p=0.1$ (grey), $1\le n\le300$. Berry--Esseen says the curves stay below $3.1$; they stay below the
dashed line at $0.4748$, as Shevtsova's constant guarantees, and the curve for $p=\frac12$ (even $n$) approaches $\frac1{\sqrt{2\pi}}\approx0.399$ (Exercise~\ref{ex:binomialsharp}); the dotted line marks that value.}
\label{fig:berryesseen}
\end{figure}

\begin{think}
The Berry--Esseen bound is uniform in $x$, so it says nothing useful about tail probabilities like $P(S_n>10\sqrt n)$, which are much smaller than $n^{-1/2}$.
Which theorem of Volume~4 would you use for those? (Chapter~7 of this volume sharpens it.)
\end{think}

\section{Several dimensions and Pearson's test}

\begin{theorem}[multivariate central limit theorem]\label{thm:multiclt}
Let $X_1,X_2,\dots$ be i.i.d.\ random vectors in $\R^d$ with $\mathbb E|X_1|^2<\infty$, mean $\mu$ and covariance matrix $\Sigma$. Then
$(S_n-n\mu)/\sqrt n\Rightarrow N(0,\Sigma)$.
\end{theorem}
\begin{proof}
Let $Y\sim N(0,\Sigma)$ (it exists: Volume~4, Theorem~6.15). For $t\in\R^d$, the variables $\langle t,X_i-\mu\rangle$ are i.i.d., centred, with variance
$t^{\mathsf T}\Sigma t$. If this is $0$ they vanish almost surely and $\langle t,(S_n-n\mu)/\sqrt n\rangle=0=\langle t,Y\rangle$ almost surely; otherwise the
one-dimensional theorem gives convergence to $N(0,t^{\mathsf T}\Sigma t)$, the law of $\langle t,Y\rangle$. Conclude by Cram\'er--Wold.
\end{proof}

The most famous application is older than the theorem. Exercise~\ref{ex:multinomial} derives from it Pearson's: if $(N_1,\dots,N_d)$ counts $n$
independent trials falling into $d$ cells with probabilities $p_1,\dots,p_d$, then
\[
\chi^2=\sum_{i=1}^d\frac{(N_i-np_i)^2}{np_i}\Rightarrow\chi^2_{d-1},
\]
the law of the sum of the squares of $d-1$ independent standard Gaussians. Large values of $\chi^2$ are evidence against the probabilities $p_i$.

\begin{example}[Weldon's dice]\label{ex:weldon}
In a letter to Francis Galton dated 2 February 1894, the zoologist Raphael Weldon reported $26\,306$ throws of twelve dice, recording each time the
number of dice showing $5$ or $6$. For fair dice this number is binomial$(12,\frac13)$. Here are the counts, with the expected values (rounded):
\begin{center}
\small
\begin{tabular}{r|rrrrr}
fives or sixes&0&1&2&3&4\\\hline
observed&185&1149&3265&5475&6114\\
expected&203&1217&3345&5576&6273
\end{tabular}
\par\medskip
\begin{tabular}{r|rrrr}
fives or sixes&5&6&7&8\\\hline
observed&5194&3067&1331&403\\
expected&5018&2927&1254&392
\end{tabular}
\par\medskip
\begin{tabular}{r|rrrr}
fives or sixes&9&10&11&12\\\hline
observed&105&14&4&0\\
expected&87&13&1&0
\end{tabular}
\end{center}
There are too few small counts and too many large ones. In 1900 Karl Pearson used these data to introduce his test, and concluded that the dice
were not fair, quoting odds of about $62\,499$ to $1$ against chance. Redone today: the last two cells expect only $1.2$ and $0.05$ throws, too few
for the $\chi^2$ approximation (a single throw in the last cell would add $20$ to the statistic), so we merge the cells $10$, $11$ and $12$. The
statistic over the eleven remaining cells is $35.5$, and a $\chi^2_{10}$ variable exceeds it with probability about $10^{-4}$: the verdict stands. The
estimated probability of a five or a six is $0.3377$ instead of $0.3333$; with that value the statistic drops to $8.2$, an ordinary value for
$\chi^2_9$ (one parameter has been estimated, which removes one degree of freedom, as Fisher explained in 1924). In 2009 Zacariah Labby rebuilt the experiment with a dice-throwing
machine, $26\,306$ throws of twelve dice again; he too found fives and sixes slightly too frequent, but traced the excess to small irregularities of
the cubes rather than to the hollowed-out pips that had long been blamed.
\end{example}

\section{The local limit theorem}

The central limit theorem speaks of intervals. For integer-valued variables one can ask for more: the probability of each single value. Volume~4,
Theorem~6.12 did it for the binomial law, with Stirling's formula. The Fourier method does it for every integer law with a finite variance,
provided the law is not trapped on a sublattice. An integer-valued $X$ has \emph{span~$1$} if the differences $x-y$ of points of its support (the
values taken with positive probability) have greatest common divisor $1$. Otherwise, with $h\ge2$ that divisor, $X$ lives in $x_0+h\Z$, and so
does $S_n-nx_0$: then $P(S_n=k)=0$ for most $k$, and the question must be asked on $nx_0+h\Z$ (Exercise~\ref{ex:lltdensity} hints at the
general picture).

\begin{lemma}\label{lem:span}
If $X$ is integer-valued with span $1$, then $|\varphi_X(t)|<1$ for $0<|t|\le\pi$.
\end{lemma}
\begin{proof}
If $|\varphi_X(t)|=1$, write $\varphi_X(t)=e^{i\theta}$. Then $\mathbb E[1-\cos(tX-\theta)]=0$, so $tX-\theta\in2\pi\Z$ almost surely, and $t(x-y)\in2\pi\Z$
for all $x,y$ in the support. By B\'ezout, $1$ is an integer combination of such differences, so $t\in2\pi\Z$, which excludes $0<|t|\le\pi$.
\end{proof}

\begin{theorem}[local limit theorem; Gnedenko, 1948]\label{thm:llt}
Let $X_1,X_2,\dots$ be i.i.d., integer-valued with span $1$, mean $\mu$ and variance $\sigma^2\in(0,\infty)$, and $\phi(x)=e^{-x^2/2}/\sqrt{2\pi}$. Then
\[
\sup_{k\in\Z}\Bigl|\sigma\sqrt n\,P(S_n=k)-\phi\Bigl(\frac{k-n\mu}{\sigma\sqrt n}\Bigr)\Bigr|\xrightarrow[n\to\infty]{}0.
\]
\end{theorem}
\begin{proof}
Let $\varphi=\varphi_{X_1}$ and $\chi(t)=e^{-i\mu t}\varphi(t)$, the characteristic function of $X_1-\mu$. Since $\varphi(t)^n=\sum_jP(S_n=j)e^{itj}$, an absolutely
convergent series, integrating term by term gives
\[
P(S_n=k)=\frac1{2\pi}\int_{-\pi}^\pi e^{-itk}\varphi(t)^n\dd t.
\]
Put $x=\frac{k-n\mu}{\sigma\sqrt n}$ and substitute $t=\frac u{\sigma\sqrt n}$; as $e^{-itk}\varphi(t)^n=e^{-iux}\chi(t)^n$,
\[
\sigma\sqrt n\,P(S_n=k)=\frac1{2\pi}\int_{|u|\le\pi\sigma\sqrt n}e^{-iux}\psi_n(u)\dd u,\qquad\psi_n(u)=\chi\Bigl(\frac u{\sigma\sqrt n}\Bigr)^n,
\]
and $\psi_n$ is the characteristic function of $(S_n-n\mu)/(\sigma\sqrt n)$. By inversion, $\phi(x)=\frac1{2\pi}\int e^{-iux}e^{-u^2/2}\dd u$. Since
$\chi(t)=1-\frac{\sigma^2t^2}2+o(t^2)$ (Volume~4, Lemma~6.6), there is $\delta\in(0,\pi]$ with $|\chi(t)|\le1-\frac{\sigma^2t^2}4\le e^{-\sigma^2t^2/4}$ for
$|t|\le\delta$; and by Lemma~\ref{lem:span} and continuity, $\eta=\max_{\delta\le|t|\le\pi}|\varphi(t)|<1$. For $A>0$ the difference in the theorem is
at most $\frac1{2\pi}$ times
\[
\begin{gathered}
\int_{|u|\le A}\bigl|\psi_n(u)-e^{-u^2/2}\bigr|\dd u+\int_{|u|>A}e^{-u^2/2}\dd u\\
+\int_{A<|u|\le\delta\sigma\sqrt n}|\psi_n(u)|\dd u+\int_{\delta\sigma\sqrt n<|u|\le\pi\sigma\sqrt n}|\psi_n(u)|\dd u,
\end{gathered}
\]
uniformly in $x$. The first term tends to $0$ by the central limit theorem (pointwise convergence of $\psi_n$, bounded by $2$ on a bounded interval).
In the third, $|\psi_n(u)|\le e^{-u^2/4}$. The fourth is at most $2\pi\sigma\sqrt n\,\eta^n\to0$. So the $\limsup$ of the difference is at most
$\frac1{2\pi}\int_{|u|>A}2e^{-u^2/4}\dd u$, which is as small as we like.
\end{proof}

\begin{figure}[ht]
\centering\fig{ch04-llt}
\caption{The exact law of $S_{30}$ for steps taking the values $0$, $1$, $3$ with probabilities $\frac12$, $\frac14$, $\frac14$ (bars), against
$\frac1{\sigma\sqrt n}\phi\bigl(\frac{k-n\mu}{\sigma\sqrt n}\bigr)$ (curve). Here $\mu=1$, $\sigma^2=\frac32$, and the support $\{0,1,3\}$ has span $1$.
The fit is already good point by point, with a visible skew that the next term of the expansion (the Edgeworth correction) would capture.}
\label{fig:llt}
\end{figure}

\section{Coming home: recurrence by Fourier analysis}

Volume~4 decided recurrence for the simple walk on $\Z^d$ by counting paths. Fourier analysis, in the hands of Kai Lai Chung and Wolfgang Fuchs (1951), decides it for every walk
with finite variance, and explains the dimension $2$ as a convergence exponent: whether $\int|t|^{-2}\dd t$ converges near $0$ in $\R^d$.

\begin{lemma}\label{lem:greenfourier}
Let $S_n$ be a walk on $\Z^d$ with i.i.d.\ steps of characteristic function $\varphi$, and $u_n=P(S_n=0)$. For $0<r<1$,
\[
\sum_{n\ge0}r^nu_n=\frac1{(2\pi)^d}\int_{[-\pi,\pi]^d}\operatorname{Re}\frac1{1-r\varphi(t)}\dd t.
\]
\end{lemma}
\begin{proof}
For $x\in\Z^d$, $\int_{[-\pi,\pi]^d}e^{i\langle t,x\rangle}\dd t$ is $(2\pi)^d$ if $x=0$ and $0$ otherwise. So by Fubini $u_n=(2\pi)^{-d}\int\mathbb E e^{i\langle t,S_n\rangle}\dd t=
(2\pi)^{-d}\int\varphi(t)^n\dd t$. Summing, $\sum_nr^n\varphi(t)^n=\frac1{1-r\varphi(t)}$, with partial sums bounded by $\frac1{1-r}$, so we may integrate term
by term. The left side is real, so only the real part of the integral survives.
\end{proof}

Note that $\operatorname{Re}\frac1z=\frac{\operatorname{Re}z}{|z|^2}>0$ when $\operatorname{Re}z>0$, which is the case for $z=1-r\varphi(t)$. We say that the
steps \emph{generate} $\Z^d$ if every point of $\Z^d$ is an integer combination of points of the support. Then the support spans $\R^d$, so it
contains a basis $v_1,\dots,v_d$ of $\R^d$; and $\varphi(t)=1$ forces $\langle t,x\rangle\in2\pi\Z$ for every $x$ in the support, hence for every $x\in\Z^d$, so
$t\in2\pi\Z^d$: in the cube $[-\pi,\pi]^d$, only $t=0$.

\begin{theorem}[recurrence on $\Z^d$]\label{thm:recurrencefourier}
Let the i.i.d.\ steps $X_i$ of a walk generate $\Z^d$.
\begin{enumerate}[label=(\alph*)]
\item If $d\ge3$, the walk is transient, whatever the law of the steps.
\item If $d\le2$, $\mathbb E|X_1|^2<\infty$ and $\mathbb EX_1=0$, the walk is recurrent.
\end{enumerate}
\end{theorem}
\begin{proof}
\emph{A lower bound.} Let $p=\min_jP(X_1=v_j)>0$ and $M=\max_j|v_j|$. For $|t|\le\frac1M$, $|\langle t,v_j\rangle|\le1$, and $1-\cos u\ge\frac{2u^2}{\pi^2}$ for
$|u|\le\pi$, so
\[
1-\operatorname{Re}\varphi(t)=\mathbb E\bigl[1-\cos\langle t,X_1\rangle\bigr]\ge\frac{2p}{\pi^2}\sum_j\langle t,v_j\rangle^2\ge c|t|^2,
\]
because $\sum_j\langle t,v_j\rangle^2$ is a positive definite quadratic form. On the compact set of $t\in[-\pi,\pi]^d$ with $|t|\ge\frac1M$, the continuous
function $1-\operatorname{Re}\varphi$ vanishes only where $\varphi(t)=1$, that is nowhere; so it is bounded below there by a positive constant, hence by
$c'|t|^2$. Shrinking $c$, $1-\operatorname{Re}\varphi(t)\ge c|t|^2$ on the whole cube.

(a) Since $|z|\ge\operatorname{Re}z$ and $1-r\operatorname{Re}\varphi=(1-r)+r(1-\operatorname{Re}\varphi)$,
\[
\operatorname{Re}\frac1{1-r\varphi(t)}\le\frac1{\operatorname{Re}(1-r\varphi(t))}\le\frac1{r\bigl(1-\operatorname{Re}\varphi(t)\bigr)}\le\frac1{rc|t|^2}.
\]
In dimension $d\ge3$, $|t|^{-2}$ is integrable near $0$ (in polar coordinates, $\int_0^1s^{-2}s^{d-1}\dd s<\infty$). So by Lemma~\ref{lem:greenfourier}, $\sum_nr^nu_n$
is bounded for $r\in[\frac12,1)$, and by monotone convergence $\sum_nu_n<\infty$: the walk is transient.

(b) Let $K=\mathbb E|X_1|^2$. From $1-\cos u\le\frac{u^2}2$ and Cauchy--Schwarz, $1-\operatorname{Re}\varphi(t)\le\frac K2|t|^2$. Since $\mathbb EX_1=0$,
$\operatorname{Im}\varphi(t)=\mathbb E[\sin\langle t,X_1\rangle-\langle t,X_1\rangle]$, and $|\sin u-u|\le u^2$ for all real $u$ (it is at most $\frac{|u|^3}6$, and at
most $1+|u|$), so $|\operatorname{Im}\varphi(t)|\le K|t|^2$. Hence for $0<|t|\le\frac1M$,
\[
\operatorname{Re}\frac1{1-\varphi(t)}=\frac{1-\operatorname{Re}\varphi(t)}{(1-\operatorname{Re}\varphi(t))^2+(\operatorname{Im}\varphi(t))^2}\ge\frac{c|t|^2}{\frac54K^2|t|^4}=\frac{4c}{5K^2|t|^2},
\]
which is \emph{not} integrable near $0$ when $d\le2$. As $r\uparrow1$, $\operatorname{Re}\frac1{1-r\varphi(t)}\to\operatorname{Re}\frac1{1-\varphi(t)}$ for every $t\ne0$,
and these functions are non-negative; by Fatou's lemma and Lemma~\ref{lem:greenfourier}, $\sum_nu_n\ge\liminf_{r\uparrow1}\sum_nr^nu_n=\infty$. The walk is
recurrent.
\end{proof}

\begin{remark}[P\'olya's constant]
For the simple walk on $\Z^3$, $\varphi(t)=\frac13(\cos t_1+\cos t_2+\cos t_3)$ is real, and Lemma~\ref{lem:greenfourier} with $r\to1$ gives
\[
\sum_{n\ge0}u_n=\frac1{(2\pi)^3}\int_{[-\pi,\pi]^3}\frac{\dd t}{1-\frac13(\cos t_1+\cos t_2+\cos t_3)}=1.5163\ldots,
\]
one of three ``triple integrals'' evaluated in closed form by G.~N.~Watson in 1939. Since $\sum_{n\ge0}u_n=\frac1{1-F}$ (Volume~4, Theorem~7.12), the
probability of ever returning is $F=0.3405\ldots$: a random walker in three dimensions comes home only about one time in three. Direct computation gives
the integral.
\end{remark}

\begin{think}
Part (b) needs a finite variance only in dimension $2$. In dimension $1$, Chapter~2 (Exercise~\ref{ex:chungfuchs}) showed that a finite \emph{mean}
$0$ is enough. Can a symmetric walk on $\Z$ be transient? The answer is yes, with heavy enough tails: see Exercise~\ref{ex:heavytail}.
\end{think}

\section{Exercises}

\begin{exercise}[$\star$]\label{ex:fellercounter}
Let $\varepsilon_1,\varepsilon_2,\dots$ be independent fair signs and $a_k=2^{-k}$. Show that $\sum_{k\le n}a_k\varepsilon_k$, normalised to have variance $1$,
converges in distribution to the uniform law on $[-\sqrt3,\sqrt3]$. Which hypothesis of Theorem~\ref{thm:feller} fails? Show on the other hand
that if $a_k=k^{-\alpha}$ with $\alpha\le\frac12$, the normalised sum is asymptotically Gaussian.
\end{exercise}

\begin{exercise}[$\star$]\label{ex:inversionsvar}
Check the formulas for $\mathbb EI_n$ and $\operatorname{Var}I_n$ in Example~\ref{ex:inversions}, and the claim that $\pi\mapsto(Y_1,\dots,Y_n)$ is injective.
\end{exercise}

\begin{exercise}[$\star$]\label{ex:chisqmean}
In Pearson's setting, show that $\mathbb E\chi^2=d-1$ for every $n$, the mean of the limit law $\chi^2_{d-1}$.
\end{exercise}

\begin{exercise}[$\star\star$]\label{ex:binomialsharp}
For the coin-tossing walk with $n$ even, show that $\liminf_n\sqrt n\sup_x|F_n(x)-\Phi(x)|\ge\frac1{\sqrt{2\pi}}\approx0.399$, improving the
constant of Proposition~\ref{prop:latticelower}. (Figure~\ref{fig:berryesseen} suggests that this is the exact limit, which is true: it follows from Esseen's
expansion for lattice laws, in Gnedenko and Kolmogorov's book.)
\end{exercise}

\begin{exercise}[$\star\star$]\label{ex:lltdensity}
(Local limit theorem for densities.) Let $X_1,X_2,\dots$ be i.i.d., centred, with variance $1$, and suppose $\int|\varphi_{X_1}|^\nu<\infty$ for some
$\nu\ge1$. Show that for $n\ge\nu$, $S_n/\sqrt n$ has a continuous density $f_n$, and that $\sup_x|f_n(x)-\phi(x)|\to0$.
\end{exercise}

\begin{exercise}[$\star\star$]\label{ex:poissonbe}
Let $N_\lambda$ be Poisson with mean $\lambda$. Show that $\sup_x\bigl|P\bigl((N_\lambda-\lambda)/\sqrt\lambda\le x\bigr)-\Phi(x)\bigr|\le3.1/\sqrt\lambda$, by
writing $N_\lambda$ as a sum of $n$ independent Poisson$(\lambda/n)$ variables and letting $n\to\infty$.
\end{exercise}

\begin{exercise}[$\star\star\star$]\label{ex:heavytail}
Let $0<\alpha<1$ and let $X$ take the value $k\in\Z\setminus\{0\}$ with probability $c|k|^{-1-\alpha}$. Show that $1-\varphi_X(t)\ge c'|t|^\alpha$ for
$|t|\le\pi$, and deduce from Lemma~\ref{lem:greenfourier} that the symmetric walk with steps distributed like $X$ is transient.
\end{exercise}

\section{Hints}

\begin{hint}{ex:fellercounter}
$\prod_{k\le n}\cos(t2^{-k})=\frac{\sin t}{2^n\sin(t2^{-n})}$: multiply by $\sin(t2^{-n})$ and use $\sin2a=2\sin a\cos a$ repeatedly. For the second part, check Lindeberg with $|X_{n,k}|\le a_1/s_n$.
\end{hint}
\begin{hint}{ex:inversionsvar}
$\sum_{k\le n}(k^2-1)=\frac{n(n+1)(2n+1)}6-n$.
\end{hint}
\begin{hint}{ex:chisqmean}
$N_i$ is binomial$(n,p_i)$.
\end{hint}
\begin{hint}{ex:binomialsharp}
Apply Theorem~\ref{thm:llt} to the Bernoulli variables $(X_i+1)/2$ to find the size of the atom of $S_n/\sqrt n$ at $0$, and redo the proof of
Proposition~\ref{prop:latticelower}.
\end{hint}
\begin{hint}{ex:lltdensity}
Copy the proof of Theorem~\ref{thm:llt}, with the integral over $|t|\le\pi$ replaced by an integral over $\R$, and use $|\varphi(t)|<1$ for $t\ne0$
(a variable with a density has no lattice) together with $|\varphi|^n\le\eta^{n-\nu}|\varphi|^\nu$ for $|t|\ge\delta$.
\end{hint}
\begin{hint}{ex:poissonbe}
For $Y$ Poisson$(m)$ with $m\le1$, bound $\mathbb E|Y-m|^3\le m^3+\mathbb EY^3$, and compute $\mathbb EY^3=m+3m^2+m^3$.
\end{hint}
\begin{hint}{ex:heavytail}
$1-\varphi_X(t)=\sum_kc|k|^{-1-\alpha}(1-\cos tk)$; keep only $1\le|k|\le\pi/|t|$ and use $1-\cos u\ge\frac{2u^2}{\pi^2}$. Then $\operatorname{Re}\frac1{1-r\varphi}\le\frac1{r(1-\varphi)}$
as $\varphi$ is real.
\end{hint}

\section{Solutions}

\begin{solution}{ex:fellercounter}
\idea A sum of independent pieces in which the first piece carries a fixed fraction of the variance cannot forget the law of that piece. Here the
pieces halve each time, and the limit is the law of a random real number written in binary.
\writeup The characteristic function of $\sum_{k\le n}2^{-k}\varepsilon_k$ is $\prod_{k\le n}\cos(t2^{-k})$. Multiplying by $\sin(t2^{-n})$ and using
$\sin2a=2\sin a\cos a$ $n$ times, $2^n\sin(t2^{-n})\prod_{k\le n}\cos(t2^{-k})=\sin t$; as $2^n\sin(t2^{-n})\to t$, the product tends to $\frac{\sin t}t$, the characteristic
function of the uniform law on $[-1,1]$ (since $\frac12\int_{-1}^1e^{itx}\dd x=\frac{\sin t}t$). By L\'evy's theorem the sums converge in law to $U$ uniform on
$[-1,1]$, of variance $\frac13$. The variances $s_n^2=\sum_{k\le n}4^{-k}\to\frac13$, so the normalised sums converge in law to $\sqrt3U$, uniform on $[-\sqrt3,\sqrt3]$.
The first summand has variance $\frac14/s_n^2\to\frac34$: condition~\eqref{eq:uan} fails, and the limit is not Gaussian.

If $a_k=k^{-\alpha}$ with $\alpha\le\frac12$, then $s_n^2=\sum_{k\le n}k^{-2\alpha}\ge\sum_{k\le n}k^{-1}\to\infty$, while $|a_k\varepsilon_k|\le1$. So $|X_{n,k}|\le1/s_n\to0$, the
indicators in $L_n(\varepsilon)$ vanish for large $n$, and Lindeberg's theorem applies.
\end{solution}

\begin{solution}{ex:inversionsvar}
\idea Independence turns everything into sums of the mean and the variance of a discrete uniform law.
\writeup The uniform law on $\{0,\dots,k-1\}$ has mean $\frac{k-1}2$ and second moment $\frac{(k-1)(2k-1)}6$, so variance $\frac{(k-1)(2k-1)}6-\frac{(k-1)^2}4=\frac{k^2-1}{12}$.
Then $\mathbb EI_n=\sum_{k\le n}\frac{k-1}2=\frac{n(n-1)}4$ and
\[
\operatorname{Var}I_n=\frac1{12}\Bigl(\frac{n(n+1)(2n+1)}6-n\Bigr)=\frac{2n^3+3n^2+n-6n}{72}=\frac{2n^3+3n^2-5n}{72}.
\]
Injectivity: $\pi(n)$ is preceded by all other values, so $Y_n$ counts the values larger than $\pi(n)$, and $\pi(n)=n-Y_n$. Removing the last position
and relabelling the remaining values increasingly by $1,\dots,n-1$ leaves a permutation of $\{1,\dots,n-1\}$ with vector $(Y_1,\dots,Y_{n-1})$; by
induction it is determined, and so is $\pi$.
\end{solution}

\begin{solution}{ex:chisqmean}
\idea The statistic is a sum of squared standardised counts; each count is binomial.
\writeup $\mathbb E(N_i-np_i)^2=np_i(1-p_i)$, so $\mathbb E\chi^2=\sum_i(1-p_i)=d-1$.
\end{solution}

\begin{solution}{ex:binomialsharp}
\idea The distribution function of a lattice law is a staircase, and a continuous function can follow it only to within half the height of each step.
The local limit theorem tells us the height of the central step.
\writeup Let $B_i=(X_i+1)/2$, Bernoulli$(\frac12)$ with span $1$, mean $\frac12$ and $\sigma=\frac12$; then $S_n=2(B_1+\dots+B_n)-n$, so $S_n=0$ means
$B_1+\dots+B_n=\frac n2$, a point where $\frac{k-n\mu}{\sigma\sqrt n}=0$. Theorem~\ref{thm:llt} gives $\frac{\sqrt n}2P(S_n=0)\to\phi(0)$. As in
Proposition~\ref{prop:latticelower}, $\sup_x|F_n(x)-\Phi(x)|\ge\frac12P(S_n=0)$, so
\[
\sqrt n\sup_x|F_n(x)-\Phi(x)|\ge\frac{\sqrt n}2P(S_n=0)\to\phi(0)=\frac1{\sqrt{2\pi}}.
\]
\end{solution}

\begin{solution}{ex:lltdensity}
\idea Inversion turns densities into integrals of characteristic functions, as in the lattice case, and the integrability of a power of $\varphi$
replaces the compactness of $[-\pi,\pi]$.
\writeup Let $\varphi=\varphi_{X_1}$. For $n\ge\nu$, $|\varphi(t/\sqrt n)|^n\le|\varphi(t/\sqrt n)|^\nu$ is integrable, so by inversion $S_n/\sqrt n$ has the continuous
density $f_n(x)=\frac1{2\pi}\int e^{-iux}\varphi(u/\sqrt n)^n\dd u$, and $\phi(x)=\frac1{2\pi}\int e^{-iux}e^{-u^2/2}\dd u$. Choose $\delta>0$ with
$|\varphi(t)|\le e^{-t^2/4}$ for $|t|\le\delta$. A variable with a density has $|\varphi(t)|<1$ for $t\ne0$ (if $|\varphi(t)|=1$ the variable would live on a
lattice $a+\frac{2\pi}t\Z$, a null set); and $\varphi(t)\to0$ as $|t|\to\infty$ by the Riemann--Lebesgue lemma. So $\eta=\sup_{|t|\ge\delta}|\varphi(t)|<1$.
Then $\sup_x|f_n(x)-\phi(x)|$ is at most $\frac1{2\pi}$ times
\[
\int_{|u|\le A}|\psi_n-e^{-u^2/2}|+\int_{|u|>A}e^{-u^2/2}+\int_{A<|u|\le\delta\sqrt n}e^{-u^2/4}+\int_{|u|>\delta\sqrt n}|\varphi(u/\sqrt n)|^n\dd u,
\]
with $\psi_n(u)=\varphi(u/\sqrt n)^n$. The first tends to $0$ by the central limit theorem, the second and third are small for large $A$, and the last is at
most $\eta^{n-\nu}\sqrt n\int|\varphi|^\nu\to0$.
\end{solution}

\begin{solution}{ex:poissonbe}
\idea The law of $N_\lambda$ does not depend on how we cut it into pieces, while the Berry--Esseen bound does; cutting into very many small pieces makes
$\rho/(\sigma^3\sqrt n)$ tend to $1/\sqrt\lambda$.
\writeup Let $m=\lambda/n\le1$ and $Y\sim$ Poisson$(m)$. Then $\sigma^2=m$ and, since $|k-m|\le k$ for $k\ge1$,
$\rho=\mathbb E|Y-m|^3\le m^3P(Y=0)+\mathbb EY^3\le m^3+(m+3m^2+m^3)\le m+5m^2$. The sum of $n$ independent copies is $N_\lambda$, so by Theorem~\ref{thm:berryesseen}
the distance is at most
\[
\frac{3.1\,\rho}{\sigma^3\sqrt n}\le\frac{3.1(m+5m^2)}{m^{3/2}\sqrt n}=\frac{3.1(1+5m)}{\sqrt{mn}}=\frac{3.1(1+5\lambda/n)}{\sqrt\lambda}.
\]
The left side does not depend on $n$; let $n\to\infty$.
\end{solution}

\begin{solution}{ex:heavytail}
\idea Transience in dimension $1$ needs $1-\varphi$ to vanish more slowly than $|t|$ at $0$, so that $\int\frac{\dd t}{1-\varphi(t)}$ converges. Heavy tails make $1-\varphi$
large: the walk makes long jumps and loses its way home.
\writeup For $0<|t|\le\pi$, keep in $1-\varphi_X(t)=\sum_{k\ne0}c|k|^{-1-\alpha}(1-\cos tk)$ only the terms with $1\le|k|\le\pi/|t|$, where $|tk|\le\pi$ and so
$1-\cos tk\ge\frac{2t^2k^2}{\pi^2}$:
\[
1-\varphi_X(t)\ge\frac{4ct^2}{\pi^2}\sum_{1\le k\le\pi/|t|}k^{1-\alpha}\ge\frac{4ct^2}{\pi^2}\int_0^{\lfloor\pi/|t|\rfloor}s^{1-\alpha}\dd s
\ge\frac{4ct^2}{\pi^2(2-\alpha)}\Bigl(\frac\pi{2|t|}\Bigr)^{2-\alpha},
\]
using $\lfloor\pi/|t|\rfloor\ge\frac\pi{2|t|}$ for $|t|\le\pi$. This is $c'|t|^\alpha$. As $X$ is symmetric, $\varphi_X$ is real and $1-r\varphi_X\ge r(1-\varphi_X)$, so
$\operatorname{Re}\frac1{1-r\varphi_X(t)}\le\frac1{rc'|t|^\alpha}$, integrable on $[-\pi,\pi]$ because $\alpha<1$. By Lemma~\ref{lem:greenfourier} and monotone convergence,
$\sum_nu_n<\infty$, and the walk is transient (Volume~4, Theorem~7.12). It has no mean ($\mathbb E|X|=\infty$ since $\alpha<1$), which is why Chapter~2's
criterion did not apply.
\end{solution}

\section*{Sources and further reading}

\begin{itemize}
\item \emph{History.} G.~P\'olya, ``\"Uber den zentralen Grenzwertsatz der Wahrscheinlichkeitsrechnung und das Momentenproblem'', \emph{Math.\ Z.}\ 8 (1920),
171--181. J.~W.~Lindeberg, ``Eine neue Herleitung des Exponentialgesetzes in der Wahrscheinlichkeitsrechnung'', \emph{Math.\ Z.}\ 15 (1922), 211--225. W.~Feller,
``\"Uber den zentralen Grenzwertsatz der Wahrscheinlichkeitsrechnung'', \emph{Math.\ Z.}\ 40 (1935), 521--559. P.~L\'evy, ``Propri\'et\'es asymptotiques des sommes de
variables al\'eatoires ind\'ependantes ou encha\^in\'ees'', \emph{J.\ Math.\ Pures Appl.}\ 14 (1935), 347--402. L.~Le Cam, ``The central limit theorem around 1935'',
\emph{Statist.\ Sci.}\ 1 (1986), 78--96. A.~C.~Berry, ``The accuracy of the Gaussian approximation to the sum of independent variates'', \emph{Trans.\ Amer.\ Math.\ Soc.}\ 49
(1941), 122--136. C.-G.~Esseen, ``On the Liapunoff limit of error in the theory of probability'', \emph{Ark.\ Mat.\ Astr.\ Fys.}\ 28A (1942), 1--19; ``A moment inequality
with an application to the central limit theorem'', \emph{Skand.\ Aktuarietidskr.}\ 39 (1956), 160--170. I.~G.~Shevtsova, ``On the absolute constants in the
Berry--Esseen type inequalities for identically distributed summands'', arXiv:1111.6554 (2011). K.~Pearson, ``On the criterion that a given system of deviations
from the probable in the case of a correlated system of variables is such that it can be reasonably supposed to have arisen from random sampling'',
\emph{Phil.\ Mag.}\ (5) 50 (1900), 157--175. R.~A.~Fisher, ``The conditions under which $\chi^2$ measures the discrepancy between observation and hypothesis'',
\emph{J.\ Roy.\ Statist.\ Soc.}\ 87 (1924), 442--450. Z.~Labby, ``Weldon's dice, automated'', \emph{Chance} 22(4) (2009), 6--13 (Weldon's data and letter, and the
modern repetition). B.~V.~Gnedenko, ``On a local limit theorem of the theory of probability'' (in Russian), \emph{Uspekhi Mat.\ Nauk} 3(3) (1948), 187--194.
K.~L.~Chung and W.~H.~J.~Fuchs, ``On the distribution of values of sums of random variables'', \emph{Mem.\ Amer.\ Math.\ Soc.}\ 6 (1951). G.~N.~Watson, ``Three triple integrals'', \emph{Quart.\ J.\ Math.}\ 10 (1939), 266--276.
\item B.~V.~Gnedenko and A.~N.~Kolmogorov, \emph{Limit Distributions for Sums of Independent Random Variables}, translated by K.~L.~Chung, Addison-Wesley, 1954. The
classical treatise: Lindeberg--Feller, local limit theorems, and the general limit laws of Chapter~5.
\item W.~Feller, \emph{An Introduction to Probability Theory and Its Applications}, vol.~II, 2nd ed., Wiley, 1971, Chapters~XV--XVI. Smoothing, Berry--Esseen and
Edgeworth expansions; our Lemmas~\ref{lem:smoothing} and~\ref{lem:cfestimate} follow it.
\item R.~Durrett, \emph{Probability: Theory and Examples}, 5th ed., Cambridge University Press, 2019, Chapter~3 (Lindeberg--Feller, Berry--Esseen,
local limit theorems) and the section on recurrence of random walks (the Chung--Fuchs criterion).
\item F.~Spitzer, \emph{Principles of Random Walk}, 2nd ed., Springer, 1976. Fourier analysis of random walks on $\Z^d$, Green functions, and P\'olya's constant.
\item D.~E.~Knuth, \emph{The Art of Computer Programming}, vol.~3, \emph{Sorting and Searching}, 2nd ed., Addison-Wesley, 1998, Section~5.1.1. Inversions and
the inversion table.
\end{itemize}
